\documentclass[11pt]{article}
\usepackage[a4paper,margin=32mm]{geometry}
\usepackage{fontspec}
\usepackage{amsmath,amssymb,mathtools}
\usepackage{enumitem}
\usepackage{microtype}
\setmainfont{DejaVu Serif}
\emergencystretch=3em
\allowdisplaybreaks
\setlist[itemize]{leftmargin=*,itemsep=0.35em}
\setlist[enumerate]{leftmargin=*,itemsep=0.35em}
\newcommand{\Fq}{\mathbb F_q}
\newcommand{\Gal}{\operatorname{Gal}}
\newcommand{\Hom}{\operatorname{Hom}}
\newcommand{\Sw}{\operatorname{Sw}}
\newcommand{\dimn}{\operatorname{dim}}
\newcommand{\eps}{\varepsilon}
\newcommand{\cO}{\mathcal O}
\newcommand{\dual}[1]{\check{#1}}
\begin{document}
\begin{center}
Bures-sur-Yvette, July 1, 1969

\vspace{1.0cm}
{\Large Letter to André Weil}
\end{center}

\vspace{0.4cm}
Dear Monsieur Weil,

Here is the determination, in the case of a function field, of the way in which the constant in the functional equation we had discussed varies. I enclose a copy of a letter from Langlands, and some calculations related to it. I have used only the statement (L) below.

\section*{1. $\lambda$-adic representations, after Serre and Grothendieck}

Let $K$ be a function field in one variable with field of constants $\Fq$, let $X$ be the set of its places, let $E$ be a number field, let $\lambda$ be a place of $E$ not dividing $q$, and let $E_\lambda$ be the corresponding local field, with ring of integers $\cO_\lambda$ and residue field $k_\lambda=\cO_\lambda/\lambda$.

A $\lambda$-adic representation of $K$ is a continuous representation of $\Gal(\overline K/K)$ in a finite-dimensional vector space over $E_\lambda$, such that there is a finite set $S$ of $X$ outside which $\sigma$ is unramified.

If $v\in X$, let $\overline F_v$ denote a corresponding Frobenius substitution, let $F_v=\overline F_v^{-1}$ be a geometric Frobenius, let $I_v$ be the inertia group of the local Galois group at $v$, and let $\sigma_v$ be the local representation defined by $\sigma$. Set
\[
 Z_v(\sigma,T)=\det(1-F_vT,\sigma_v^{I_v})^{-1}
\]
(the action of Frobenius on inertia invariants), and
\[
 Z(\sigma,T)=\prod_v Z_v(\sigma,T)\in E_\lambda[[T]].
\]

\paragraph{Example.} A $\lambda$-adic representation of rank $1$ of $K$ is the same as a homomorphism from the idèle class group of $K$ to $\cO_\lambda^*$, continuous for the discrete topology of $\cO_\lambda^*$. We denote by $\mathbf 1(n)$ the representation corresponding to $x\mapsto |x|^{-n}$, $n\in\mathbb Z$, and put
\[
 \sigma(n)=\sigma\otimes \mathbf 1(n).
\]

A theorem of Grothendieck asserts that, if $\sigma$ is a $\lambda$-adic representation of $K$, then:
\begin{enumerate}[label=\alph*)]
\item $Z(\sigma,T)$ is a rational function, in $E_\lambda(T)$;
\item it satisfies a functional equation
\[
 Z(\dual\sigma(1),T)=\eps(\sigma,T) Z(\sigma,T),
\]
where $\eps(\sigma,T)$ is a monomial. We put
\[
 \eps(\sigma,T)=q^{(g-1)\dimn\sigma}\,\eps_0(\sigma,T),
 \qquad \eps_0(\sigma)=\eps_0(\sigma,1),
\]
the latter number being a $\lambda$-adic unit;
\item if $\sigma$ also denotes the $\lambda$-adic sheaf defined by $\sigma$, the direct image of its restriction to the generic point, then one also has
\[
 \eps(\sigma,T)=\det(-FT,H^*(\sigma))
 =\prod_i \det(-FT,H^i(\sigma))^{(-1)^i}.
\]
\end{enumerate}

If $S$ is a finite subset of $X$, put
\[
 Z_S(\sigma,T)=\prod_{v\notin S} Z_v(\sigma,T)
\]
and
\begin{align*}
 \eps_S(\sigma,T)
 &=\eps(\sigma,T)\prod_{v\in S}\det(-F_vT,\sigma_v^{I_v})^{-1} \\
 &=\text{the monomial asymptotic to }Z_S(\sigma,T)\text{ as }T\to\infty \\
 &=\det(-FT,H^*(\sigma_S)),
\end{align*}
where $\sigma_S$ is the $\lambda$-adic sheaf $\sigma$, restricted to $X-S$ and extended by $0$ at the $v\in S$. We put similarly
\[
 \eps_S(\sigma,T)=q^{(g-1)\dimn\sigma}\,\eps_{0,S}(\sigma,T).
\]

The conductors are defined by
\[
 m_v(\sigma_S)=
 \begin{cases}
 \dimn(\sigma)-\dimn(\sigma_v^{I_v})+
 \dimn \Hom_{I_v}(\Sw_v,\sigma_v),& v\notin S,\\[0.4em]
 \dimn(\sigma)+\dimn \Hom_{I_v}(\Sw_v,\sigma_v),& v\in S,
 \end{cases}
\]
where, for instance, one sets
\[
 \dimn\Hom_{I_v}(\Sw_v,\sigma_v)=
 \dimn\Hom_{\overline I_v}(\Sw_v,\sigma_v\bmod \lambda).
\]
Here $\Sw_v$ is the Swan representation of a sufficiently large quotient $I'_v$ of $I_v$; it is a projective $\cO_\lambda[I'_v]$-module. The mention of $S$ is omitted when $S=\varnothing$.

Let now $L$ be a nonempty set of places of $E$, prime to $q$, and let $S$ be a finite subset of $X$. A compatible system of $\lambda$-adic representations, $\lambda\in L$, with exceptional set $S$, consists, for each $\lambda\in L$, of a $\lambda$-adic representation $\sigma_\lambda$ of $K$, satisfying:
\begin{enumerate}[label=(\roman*)]
\item each $\sigma_\lambda$ is unramified outside $S$;
\item for $v\notin S$, $Z_v(\sigma_\lambda,T)$ lies in $E(T)$ and is independent of $\lambda\in L$.
\end{enumerate}
Then put
\[
 Z_v(\sigma,T)=Z_v(\sigma_\lambda,T)\quad(v\notin S),
 \qquad
 Z_S(\sigma,T)=\prod_{v\notin S} Z_v(\sigma,T),
\]
and
\[
 \eps_S(\sigma,T)=\eps_S(\sigma_\lambda,T),
\]
these quantities being independent of $\lambda$; indeed, $\eps_S$ can be defined in terms of $Z_S$.

\section*{2. The statement (L)}

We shall admit the following statement. Let $\sigma$ and $\rho$ be two complex representations of a finite quotient of $\Gal(\overline K/K)$. Suppose that their conductors are disjoint. With the preceding notation, one has
\begin{equation}\tag{L}
\eps_0(\sigma\otimes\rho,T)
=\eps_0(\sigma,T)^{\dimn\rho}\,
 \eps_0(\rho,T)^{\dimn\sigma}
 \prod_{v\in X}\det\sigma(F_v)^{-m_v(\rho)}
 \prod_{v\in X}\det\rho(F_v)^{-m_v(\sigma)}.
\end{equation}

\section*{3. Variation of the constants}

\paragraph{Theorem.} Suppose that $L$ is infinite, and let $\sigma$ and $\rho$ be two compatible systems of $\lambda$-adic representations, $\lambda\in L$, with disjoint exceptional sets. Then, for every $\lambda\in L$,
\begin{equation}\tag{*}
\eps_0(\sigma_\lambda\otimes\rho_\lambda,T)
=\eps_0(\sigma_\lambda,T)^{\dimn\rho}\,
 \eps_0(\rho_\lambda,T)^{\dimn\sigma}
 \prod_{v\in X}\det\sigma_\lambda(F_v)^{-m_v(\rho_\lambda)}
 \prod_{v\in X}\det\rho_\lambda(F_v)^{-m_v(\sigma_\lambda)};
\end{equation}
thus the formula (L) remains valid.

Let $S$ and $T$ be disjoint exceptional sets for $\sigma$ and $\rho$.

\paragraph{Lemma 1.} For a fixed $\lambda$, the formula $(*)$ is equivalent to
\begin{equation}\tag{$*_{S,T}$}
\begin{aligned}
\eps_{0,S\cup T}(\sigma_\lambda\otimes\rho_\lambda,T)
&=\eps_{0,S}(\sigma_\lambda,T)^{\dimn\rho}\,
  \eps_{0,T}(\rho_\lambda,T)^{\dimn\sigma} \\
&\quad\cdot \prod_{v\in T}\det\sigma_\lambda(F_v)^{-m_v(\rho_\lambda)}
 \prod_{v\in S}\det\rho_\lambda(F_v)^{-m_v(\sigma_\lambda,S)} .
\end{aligned}
\end{equation}

If $(A)$, respectively $(B)$, denotes the ratio of the two sides of $(*)$, respectively of $(*_{S,T})$, then the quotient $(A)/(B)$ appears as a product of local terms, for $v\in S\cup T$, all of which are checked to be $1$. If $v\in S$, one has to check
\[
 \det(-F_vT,\sigma_v^{I_v}\otimes\rho_v)
 =\det(-F_vT,\sigma_v^{I_v})^{\dimn\rho}
   \det\rho(F_v)^{\dimn\sigma_v^{I_v}},
\]
which is evident.

If $\sigma$ is a representation of $\Gal(\overline K/K)$ in a finite-dimensional $k_\lambda$-vector space, define $Z(\sigma,T)$ and $Z_S(\sigma,T)$ in $k_\lambda[[T]]$ by the same formulas as before. It is trivial by reduction modulo $\lambda$ that, if $\sigma$ is a representation of $\Gal(\overline K/K)$ in a finite free $\cO_\lambda$-module, with conductor contained in $S$, and if $\sigma/\lambda$ is the $k_\lambda$-representation obtained by reduction modulo $\lambda$, then:

\paragraph{Lemma 2.}
\begin{enumerate}[label=(\roman*)]
\item $Z_S(\sigma/\lambda,T)$ is a rational function. We shall denote by
\[
 q^{(g-1)\dimn\sigma}\eps_0(\sigma/\lambda,T)
\]
the monomial asymptotic to $Z_S(\sigma/\lambda,T)$ as $T\to\infty$.

\item If $\sigma$ and $\rho$ are as above, with conductors contained in disjoint finite subsets $S$ and $T$ of $X$, the following conditions are equivalent:
\begin{align*}
\text{(a)}\quad
\eps_{0,S\cup T}(\sigma\otimes\rho,T)
&=\eps_{0,S}(\sigma,T)^{\dimn\rho}\,
  \eps_{0,T}(\rho,T)^{\dimn\sigma}\\
&\quad\cdot
\prod_{v\in T}\det\sigma(F_v)^{-m_v(\rho)}
\prod_{v\in S}\det\rho(F_v)^{-m_v(\sigma_S)},\\[0.5em]
\text{(b)}\quad
\eps_{0,S\cup T}(\sigma/\lambda\otimes\rho/\lambda,T)
&=\eps_{0,S}(\sigma/\lambda,T)^{\dimn\rho}\,
  \eps_{0,T}(\rho/\lambda,T)^{\dimn\sigma}\\
&\quad\cdot
\prod_{v\in T}\det(\sigma/\lambda)(F_v)^{-m_v(\rho/\lambda)}\\
&\quad\cdot
\prod_{v\in S}\det(\rho/\lambda)(F_v)^{-m_v((\sigma/\lambda)_S)} .
\end{align*}
\end{enumerate}

We shall deduce from (L) that:

\paragraph{Corollary 3.}
\begin{enumerate}[label=(\roman*)]
\item If $\sigma$ is a representation of $\Gal(\overline K/K)$ in a finite-dimensional $k_\lambda$-vector space, the functions $Z_S(\sigma,T)\in k_\lambda[[T]]$ are rational.
\item If $\sigma$ and $\rho$ are two such representations, with conductors contained in disjoint finite sets $S$ and $T$, formula (b) is always valid.
\end{enumerate}

To prove (i), note that, after enlarging $E$ and $k_\lambda$ if necessary, Brauer's theorem tells us that $\sigma$ lifts virtually to a representation of $\Gal(\overline K/K)$ in a finite-dimensional vector space over $E$; one then applies 2.(c) to this representation. Assertion (ii) is proved in the same way, starting from (L) and Lemma 1. It follows that assertion 2.(ii)(a) is always true, and hence that $(*_{S,T})$ is true modulo $\lambda$.

The same technique of reduction to the case of ordinary $L$-series, or another method, makes it possible to verify that
\[
 \deg_T\bigl(\eps_S(\sigma_\lambda,T)\bigr)
 =\dimn(\sigma)(2-2g)-\sum_v m_v(\sigma_\lambda,S)
\]
(the Ogg--Shafarevich--Grothendieck formula). The left-hand side of this formula is independent of $\lambda\in L$, so that $\sum_v m_v(\sigma_\lambda,S)$ is independent of $\lambda\in L$. Hence there exists an infinite subset $L_1$ of $L$ such that the integers $m_v(\sigma_\lambda,S)$, for $v\in X$ and $\lambda\in L_1$, are independent of $\lambda$.

The two sides of formula $(*_{S,T})$ are then independent of $\lambda\in L_1$; since they are equal modulo $\lambda$ for infinitely many $\lambda$, they are equal for $\lambda\in L_1$. To finish the proof, it remains to check:

\paragraph{Lemma 4.} Suppose that $L$ is infinite, and let $\sigma$ be a compatible system of $\lambda$-adic representations, with exceptional set $S$. Then the integers $m_v(\sigma_\lambda,S)$ are independent of $\lambda\in L$.

Apply the preceding results to $\sigma$ and to a finite-order character $\chi$ of the idèle class group of $K$, with values in a finite extension of $E$, and with conductor $T$ disjoint from $S$. We already know that there are integers $m_v$ -- namely the $m_v(\sigma_\lambda,S)$ for $\lambda\in L_1$ -- such that
\[
\begin{aligned}
 \eps_{0,S\cup T}(\sigma\otimes\chi,T)
 &=\eps_{0,S}(\sigma,T)\eps_{0,T}(\chi,T)^{\dimn\sigma}
 \prod_v \det\sigma(F_v)^{-m_v(\chi_T)}\\
 &\quad\cdot \prod_v \chi(F_v)^{-m_v},
\end{aligned}
\]
and, on the other hand,
\[
\begin{aligned}
 \eps_{0,S\cup T}(\sigma\otimes\chi,T)
 &=\eps_{0,S}(\sigma,T)\eps_{0,T}(\chi,T)^{\dimn\sigma}
 \prod_v \det\sigma(F_v)^{-m_v(\chi_T)}\\
 &\quad\cdot \prod_v \chi(F_v)^{-m_v(\sigma_\lambda,S)}
 \pmod{\lambda},
\end{aligned}
\]
so that
\[
 \prod_{v\in S}\chi(F_v)^{m_v-m_v(\sigma_\lambda,S)}=1\pmod{\lambda}.
\]
Since this is true for every $\chi$, it implies that $m_v(\sigma_\lambda,S)=m_v$, which completes the proof.

\vspace{1cm}
\begin{center}
Yours,

\vspace{0.5cm}
P. Deligne
\end{center}
\end{document}
